\documentclass[a4paper,10pt,brazil]{article}
\usepackage{provastyle}
\usepackage{menukeys}

\begin{document}
\makeexamtitlepage{UFOP}{BCC222: Programação Funcional}{José Romildo}{2026/1}{Primeira Prova}{02}
\begin{multicols}{2}
\questionTitle{1}{Garantias da função pura}
\begin{flushright}
  \footnotesize
  \textit{\textbf{1: Paradigmas de programação}}\\
  \textit{c01-transparencia-referencial-03}\\

\end{flushright}
Uma \textbf{função pura} é a base fundamental da construção de software no paradigma declarativo funcional. Assinale a alternativa que descreve corretamente as duas garantias absolutas oferecidas por uma função pura.

\begin{enumerate}
  \item 1) Ela possui uma assinatura de tipo polimórfica explícita. 2) Ela substitui automaticamente operações de ponto flutuante por inteiros para evitar imprecisões matemáticas.

  \item 1) Ela sempre avalia para o mesmo resultado quando recebe os mesmos argumentos. 2) Ela não causa nenhum efeito colateral observável fora de seu escopo.

  \item 1) Ela não utiliza declarações de variáveis locais em seu corpo. 2) Ela só pode receber matrizes e listas imutáveis como parâmetros.

  \item 1) Ela avalia todos os seus argumentos de forma estrita (\emph{strict evaluation}) antes de executar. 2) Ela não depende de bibliotecas de terceiros.

  \item 1) Ela é garantida de terminar (não entrar em loop infinito). 2) Ela executa obrigatoriamente em tempo constante $\mathcal{O}(1)$.


\end{enumerate}

\questionTitle{2}{Assinatura paramétrica sem restrições}
\begin{flushright}
  \footnotesize
  \textit{\textbf{7: Polimorfismo}}\\
  \textit{cap07-param-assinatura-pura}\\

\end{flushright}
Qual das assinaturas abaixo expressa uma função com polimorfismo paramétrico puro, isto é, sem restrições de classe de tipo sobre a variável \haskellinline{a}?
\begin{enumerate}
  \item \mintinline{haskell}{(+) :: Num a => a -> a -> a}
  \item \mintinline{haskell}{show :: Show a => a -> String}
  \item \mintinline{haskell}{ord :: Char -> Int}
  \item \mintinline{haskell}{reverse :: [a] -> [a]}
  \item \mintinline{haskell}{sqrt :: Floating a => a -> a}

\end{enumerate}

\questionTitle{3}{Recursão múltipla em Fibonacci}
\begin{flushright}
  \footnotesize
  \textit{\textbf{9: Funções recursivas}}\\
  \textit{c09-formas-fibonacci-02}\\

\end{flushright}
Na definição
\begin{haskellcode}
fib n
  | n == 0 = 0
  | n == 1 = 1
  | n > 1  = fib (n - 2) + fib (n - 1)
\end{haskellcode}
por que dizemos que há recursão múltipla?

\begin{enumerate}
  \item Porque a função usa duas guardas antes do caso recursivo.

  \item Porque o caso recursivo realiza mais de uma chamada à própria função.

  \item Porque a função chama outra função auxiliar com acumuladores.

  \item Porque a função retorna uma lista com vários elementos.

  \item Porque há mais de um caso base.


\end{enumerate}

\questionTitle{4}{Equivalência entre if e guardas}
\begin{flushright}
  \footnotesize
  \textit{\textbf{5: Expressão condicional}}\\
  \textit{c05-if-vs-guardas-03}\\

\end{flushright}
Qual definição com guardas é equivalente à função abaixo?
\begin{haskellcode}
max2 x y = if x > y then x else y
\end{haskellcode}

\begin{enumerate}
  \item \begin{haskellcode}
max2 x y
  | x < y     = x
  | otherwise = y
\end{haskellcode}

  \item \begin{haskellcode}
max2 x y
  | x >= y    = y
  | otherwise = x
\end{haskellcode}

  \item \begin{haskellcode}
max2 x y
  | x > y     = x
  | otherwise = y
\end{haskellcode}

  \item \begin{haskellcode}
max2 x y = if x > y then y else x
\end{haskellcode}

  \item \begin{haskellcode}
max2 x y
  | x > y     = True
  | otherwise = False
\end{haskellcode}


\end{enumerate}

\questionTitle{5}{Identificação do caso base em leitura com sentinela}
\begin{flushright}
  \footnotesize
  \textit{\textbf{10: Ações de E/S recursivas}}\\
  \textit{c10-lacos-io-02}\\

\end{flushright}
Considere a ação:

\begin{haskellcode}
lerESomar :: IO Integer
lerESomar = do
  n <- readLn
  if n == 0
    then return 0
    else do
      somaResto <- lerESomar
      return (n + somaResto)
\end{haskellcode}

Qual trecho corresponde ao caso base da recursão?

\begin{enumerate}
  \item \haskellinline|n <- readLn|, pois toda leitura encerra uma iteração.

  \item O código não possui caso base, pois a ação \haskellinline|lerESomar| aparece no próprio corpo.

  \item \haskellinline|somaResto <- lerESomar|, pois a chamada recursiva define quando a repetição termina.

  \item \haskellinline|n == 0| junto com \haskellinline|return 0|, pois o sentinela indica que a leitura deve parar e produz a soma neutra.

  \item \haskellinline|return (n + somaResto)|, pois toda expressão com \haskellinline|return| é necessariamente o caso base.


\end{enumerate}

\questionTitle{6}{Identificadores alfanuméricos válidos}
\begin{flushright}
  \footnotesize
  \textit{\textbf{3: Expressões e definições}}\\
  \textit{c03-arquivos-identificadores-03}\\

\end{flushright}
De acordo com as regras para identificadores alfanuméricos usados como nomes de variáveis ou funções, qual alternativa apresenta um identificador válido?

\begin{enumerate}
  \item \haskellinline|2valor|.

  \item \haskellinline|valor_final'|.

  \item \haskellinline|ValorFinal|.

  \item \haskellinline|valor-final|.

  \item \haskellinline|module|.


\end{enumerate}

\questionTitle{7}{Chamada recursiva com acumulador}
\begin{flushright}
  \footnotesize
  \textit{\textbf{10: Ações de E/S recursivas}}\\
  \textit{c10-acumuladores-cauda-02}\\

\end{flushright}
Considere a versão com acumulador:

\begin{haskellcode}
lerESomar :: Integer -> IO Integer
lerESomar total = do
  n <- readLn
  if n == 0
    then return total
    else lerESomar (total + n)
\end{haskellcode}

Qual afirmação é a mais precisa?

\begin{enumerate}
  \item A chamada recursiva deixa de ser recursiva porque o acumulador substitui a função original.

  \item O uso de acumulador garante, por si só, uso constante de memória em qualquer implementação Haskell.

  \item A função soma todos os valores apenas no caso base, sem atualizar o acumulador durante a recursão.

  \item A função passa a ser pura, pois o acumulador elimina o tipo \haskellinline|IO|.

  \item No ramo recursivo, a chamada \haskellinline|lerESomar (total + n)| é a última ação executada naquele ramo.


\end{enumerate}

\questionTitle{8}{Conversão de operador infixo para função prefixa}
\begin{flushright}
  \footnotesize
  \textit{\textbf{3: Expressões e definições}}\\
  \textit{c03-aplicacao-notacoes-03}\\

\end{flushright}
Qual alternativa é uma reescrita prefixa equivalente à expressão \haskellinline|100 `mod` (5 * 2)|?

\begin{enumerate}
  \item \haskellinline|mod (100, 5 * 2)|.

  \item \haskellinline|mod 100 (5 * 2)|.

  \item \haskellinline|(mod 100 5) * 2|.

  \item \haskellinline|100 mod (5 * 2)|.

  \item \haskellinline|mod (100 * 5) 2|.


\end{enumerate}

\questionTitle{9}{Rastreamento da potência de dois}
\begin{flushright}
  \footnotesize
  \textit{\textbf{9: Funções recursivas}}\\
  \textit{c09-rastreamento-potdois-01}\\

\end{flushright}
Considere:
\begin{haskellcode}
potDois :: Integer -> Integer
potDois n
  | n == 0 = 1
  | n > 0  = 2 * potDois (n - 1)
\end{haskellcode}
Qual é o valor de \haskellinline|potDois 3|?

\begin{enumerate}
  \item \haskellinline|6|

  \item \haskellinline|8|

  \item \haskellinline|4|

  \item \haskellinline|9|

  \item \haskellinline|16|


\end{enumerate}

\questionTitle{10}{Tipos de ações de entrada}
\begin{flushright}
  \footnotesize
  \textit{\textbf{8: Programas interativos}}\\
  \textit{c08-tipos-acoes-02}\\

\end{flushright}
Qual é a assinatura de tipo correta de \haskellinline|readLn| e o que ela representa?

\begin{enumerate}
  \item \haskellinline|IO ()| — uma ação que executa a entrada, mas descarta qualquer valor produzido.

  \item \haskellinline|Read a => IO a| — uma ação de E/S que, quando executada, lê uma linha e tenta convertê-la para um valor de tipo lido pela classe Read.

  \item \haskellinline|Read a => a| — um valor puro obtido diretamente, sem necessidade de \haskellinline|IO|.

  \item \haskellinline|String -> Read a => a| — uma função pura que recebe uma string e retorna o valor lido.

  \item \haskellinline|() -> Read a => IO a| — uma função que recebe o valor unitário e então constrói a ação.


\end{enumerate}

\questionTitle{11}{Assinaturas de operadores sobrecarregados}
\begin{flushright}
  \footnotesize
  \textit{\textbf{7: Polimorfismo}}\\
  \textit{cap07-adhoc-operador-param}\\

\end{flushright}
Considere o operador ou função \haskellinline|div|. Qual alternativa descreve corretamente sua assinatura de tipo mais geral e o tipo de restrição envolvida?
\begin{enumerate}
  \item A operação é resolvida por herança de objetos em tempo de execução, não pelo sistema de classes de tipo.
  \item Assinatura: \haskellinline|a -> a -> a|. A operação é paramétrica pura, sem restrições.
  \item Assinatura: \haskellinline|Integral a => a -> a -> a|. A operação é sobrecarregada por uma classe de tipo.
  \item A operação não possui assinatura polimórfica; cada tipo numérico usa uma função com nome diferente.
  \item Assinatura: \haskellinline|Double -> Double -> Double|. A operação exige que os argumentos sejam sempre \haskellinline{Double}.

\end{enumerate}

\questionTitle{12}{Explorando módulos com :browse}
\begin{flushright}
  \footnotesize
  \textit{\textbf{2: Ambiente de desenvolvimento}}\\
  \textit{c02-ghci-modulos-03}\\

\end{flushright}
Após adicionar um módulo à sessão atual do GHCi, o desenvolvedor digita o comando \texttt{:browse Data.List}. Qual será o resultado visível no terminal?

\begin{enumerate}
  \item O GHCi procurará erros de compilação (\emph{warnings}) ignorados nas funções internas daquele módulo específico.

  \item O GHCi listará as assinaturas de tipo de todas as funções e definições de dados que estão exportadas e disponíveis publicamente naquele módulo.

  \item O GHCi exibirá o código-fonte C gerado nos bastidores para cada função presente em \texttt{Data.List}.

  \item O GHCi abrirá uma nova aba no navegador web (\emph{browser}) do sistema operacional exibindo a documentação HTML oficial do módulo.

  \item O GHCi removerá o módulo do escopo ativo, fechando as portas para futuras avaliações matemáticas.


\end{enumerate}

\questionTitle{13}{Próxima equação da função}
\begin{flushright}
  \footnotesize
  \textit{\textbf{5: Expressão condicional}}\\
  \textit{c05-guardas-exaustividade-03}\\

\end{flushright}
Considere a definição:
\begin{haskellcode}
h x
  | x > 0 = "Positivo"
h x = "Não positivo"
\end{haskellcode}
Qual é o resultado de \haskellinline|h 0|?

\begin{enumerate}
  \item Ocorre um erro, pois a primeira equação não tem \haskellinline|otherwise|.

  \item \haskellinline|"Não positivo"|

  \item \haskellinline|0|

  \item O compilador rejeita a definição por possuir duas equações para a mesma função.

  \item \haskellinline|"Positivo"|


\end{enumerate}

\questionTitle{14}{Compatibilidade dos ramos}
\begin{flushright}
  \footnotesize
  \textit{\textbf{5: Expressão condicional}}\\
  \textit{c05-if-regras-03}\\

\end{flushright}
Qual expressão abaixo é rejeitada porque os ramos \haskellinline|then| e \haskellinline|else| não produzem valores do mesmo tipo?

\begin{enumerate}
  \item \haskellinline|if odd 5 then 1 else 0|

  \item \haskellinline|if odd 5 then "ímpar" else 0|

  \item \haskellinline|if odd 5 then "ímpar" else "par"|

  \item \haskellinline|if odd 5 then [1,2] else [3,4]|

  \item \haskellinline|if odd 5 then 'i' else 'p'|


\end{enumerate}

\questionTitle{15}{Guardas sobrepostas}
\begin{flushright}
  \footnotesize
  \textit{\textbf{5: Expressão condicional}}\\
  \textit{c05-guardas-sintaxe-semantica-03}\\

\end{flushright}
Considere a função:
\begin{haskellcode}
f x
  | x > 0     = 1
  | x > 10    = 2
  | otherwise = 3
\end{haskellcode}
Qual é o valor de \haskellinline|f 20|?

\begin{enumerate}
  \item Ocorre um erro, pois duas guardas são verdadeiras.

  \item O compilador rejeita a função, pois as guardas se sobrepõem.

  \item \haskellinline|2|

  \item \haskellinline|1|

  \item \haskellinline|3|


\end{enumerate}

\questionTitle{16}{Estratégias de Redução e Confluência}
\begin{flushright}
  \footnotesize
  \textit{\textbf{1: Paradigmas de programação}}\\
  \textit{c01-mecanismo-reducao-03}\\

\end{flushright}
Ao avaliar uma expressão matemática em Haskell, o sistema pode teoricamente adotar diferentes caminhos de avaliação, como a redução \emph{de dentro para fora} (aplicativa) ou \emph{de fora para dentro} (normal/preguiçosa). 

Graças à pureza e à ausência de efeitos colaterais na linguagem, qual garantia teórica fundamental o programador possui independentemente do caminho escolhido pelo compilador?

\begin{enumerate}
  \item A garantia de que a execução da expressão consumirá exatamente a mesma quantidade de milissegundos na CPU do usuário.

  \item A garantia de que, caso ambas as estratégias cheguem a um resultado final sem entrar em \emph{loop} infinito, este resultado será sempre o mesmo (Teorema de Church-Rosser).

  \item A garantia de que os tipos numéricos fracionários estarão imunes a qualquer tipo de erro de precisão de ponto flutuante do padrão IEEE 754.

  \item A garantia de que a quantidade de memória RAM alocada dinamicamente durante o processo de avaliação será anulada pelo \emph{Garbage Collector}.

  \item A garantia de que as funções envolvidas na expressão serão automaticamente vetorizadas e enviadas para execução na GPU.


\end{enumerate}

\questionTitle{17}{Equivalência com guardas}
\begin{flushright}
  \footnotesize
  \textit{\textbf{5: Expressão condicional}}\\
  \textit{c05-if-aninhamento-03}\\

\end{flushright}
Qual definição com guardas é funcionalmente equivalente à função abaixo?
\begin{haskellcode}
sinal n = if n < 0 then -1 else if n > 0 then 1 else 0
\end{haskellcode}

\begin{enumerate}
  \item \begin{haskellcode}
sinal n
  | n <= 0    = -1
  | otherwise = 1
\end{haskellcode}

  \item \begin{haskellcode}
sinal n
  | n < 0     = -1
  | otherwise = 1
\end{haskellcode}

  \item \begin{haskellcode}
sinal n
  | n < 0     = -1
  | n > 0     = 1
  | otherwise = 0
\end{haskellcode}

  \item \begin{haskellcode}
sinal n
  | n > 0     = -1
  | n < 0     = 1
  | otherwise = 0
\end{haskellcode}

  \item \begin{haskellcode}
sinal n = if n == 0 then -1 else 1
\end{haskellcode}


\end{enumerate}

\questionTitle{18}{Relação entre Ord e Eq}
\begin{flushright}
  \footnotesize
  \textit{\textbf{7: Polimorfismo}}\\
  \textit{cap07-hierarquia-ord-implica-eq}\\

\end{flushright}
Em Haskell, \haskellinline{Ord} é uma subclasse de \haskellinline{Eq}. Qual consequência prática decorre disso?
\begin{enumerate}
  \item A relação vale apenas para tipos numéricos, como \haskellinline{Int} e \haskellinline{Double}.
  \item \haskellinline{Ord} permite imprimir valores, enquanto \haskellinline{Eq} permite convertê-los com \haskellinline{read}.
  \item Todo tipo que é instância de \haskellinline{Eq} automaticamente se torna instância de \haskellinline{Ord}.
  \item \haskellinline{Ord} e \haskellinline{Eq} são sinônimos; implementar uma delas implementa a outra sem qualquer código adicional.
  \item Para que um tipo seja instância de \haskellinline{Ord}, ele também deve ser instância de \haskellinline{Eq}.

\end{enumerate}

\questionTitle{19}{Ação recursiva sem caso base não produz lista parcial}
\begin{flushright}
  \footnotesize
  \textit{\textbf{10: Ações de E/S recursivas}}\\
  \textit{c10-bordas-robustez-02}\\

\end{flushright}
Considere a ação:

\begin{haskellcode}
leParaSempre :: IO [Integer]
leParaSempre = do
  x <- readLn
  xs <- leParaSempre
  return (x:xs)
\end{haskellcode}

Qual alternativa descreve melhor o problema dessa definição como ação de E/S?

\begin{enumerate}
  \item Ela falha por erro de tipo, pois \haskellinline|x:xs| nunca pode ser usado dentro de \haskellinline|IO|.

  \item Ela imprime cada número lido automaticamente, mesmo sem chamada a \haskellinline|print|.

  \item Como não há caso base, a ação precisa continuar executando leituras indefinidamente antes de produzir o valor \haskellinline|IO [Integer]| esperado.

  \item Ela é segura porque listas em Haskell são preguiçosas, então o programa pode usar o primeiro elemento antes de ler os demais.

  \item Ela sempre produz a lista vazia, pois \haskellinline|return| ignora todos os valores lidos.


\end{enumerate}

\questionTitle{20}{A sintaxe de múltiplas linhas}
\begin{flushright}
  \footnotesize
  \textit{\textbf{2: Ambiente de desenvolvimento}}\\
  \textit{c02-fluxo-trabalho-03}\\

\end{flushright}
Ao tentar copiar uma definição funcional de 4 linhas de um PDF e colar diretamente no \textbf{prompt} do GHCi, o terminal costuma gerar erros acusando que identificadores não estão em escopo, pois avalia cada linha independentemente. 

Qual é a solução mais segura, utilizando delimitadores, para que o GHCi aceite todo o bloco colado como uma única construção sintática?

\begin{enumerate}
  \item Adicionar barras duplas \haskellinline{//} no fim de cada linha para avisar o interpretador de que o comando continua na próxima.

  \item Utilizar aspas triplas \texttt{"""} tanto no topo quanto na base do trecho copiado.

  \item Involucrar o bloco copiado digitando \texttt{:\{} antes de colar, e \texttt{:\}} ao finalizar.

  \item Inserir o comando de avaliação em lote \texttt{:batch} na linha imediatamente anterior.

  \item Envolver a expressão utilizando parênteses duplos \texttt{((} no início e \texttt{))} no fim.


\end{enumerate}

\questionTitle{21}{Tipo produzido por return no caso base}
\begin{flushright}
  \footnotesize
  \textit{\textbf{10: Ações de E/S recursivas}}\\
  \textit{c10-return-tipos-02}\\

\end{flushright}
Analise a ação recursiva:

\begin{haskellcode}
lerLista :: IO [Int]
lerLista = do
  x <- readLn
  if x == 0
    then return []
    else do
      resto <- lerLista
      return (x:resto)
\end{haskellcode}

Tecnicamente, o que a expressão \haskellinline|return []| realiza no caso base?

\begin{enumerate}
  \item Ela cria uma ação de tipo \haskellinline|IO [Int]| que produz a lista vazia quando executada.

  \item Ela sinaliza ao compilador que a recursão deve ser transformada em um laço imperativo.

  \item Ela lê a próxima lista da entrada padrão, usando \haskellinline|readLn| implicitamente.

  \item Ela encerra todo o programa e descarta as chamadas recursivas anteriores.

  \item Ela imprime \haskellinline|[]| na tela e retorna \haskellinline|()|.


\end{enumerate}

\questionTitle{22}{Escopo e ordem da declaração default}
\begin{flushright}
  \footnotesize
  \textit{\textbf{7: Polimorfismo}}\\
  \textit{cap07-defaulting-declaracao-default}\\

\end{flushright}
Considere a declaração:

\begin{minted}{haskell}
default (Int, Integer, Float)
\end{minted}

Qual afirmação é correta?

\begin{enumerate}
  \item Ela converte todos os literais do programa para \haskellinline{Int}, mesmo quando o contexto exige \haskellinline{Double}.
  \item Ela vale globalmente para todos os módulos importados e para todos os programas compilados depois.
  \item Ela transforma \haskellinline{Int} em subclasse de \haskellinline{Float}.
  \item Ela permite resolver ambiguidades envolvendo apenas \haskellinline{Read}, mesmo sem restrição numérica ou regras estendidas.
  \item Ela altera, no módulo em que aparece, a ordem dos candidatos usados para resolver certas ambiguidades numéricas.

\end{enumerate}

\questionTitle{23}{Correção de divisão envolvendo length}
\begin{flushright}
  \footnotesize
  \textit{\textbf{7: Polimorfismo}}\\
  \textit{cap07-conversao-length-divisao}\\

\end{flushright}
Por que a expressão abaixo é problemática?

\begin{minted}{haskell}
length xs / 2
\end{minted}

Qual reescrita é a forma idiomática para obter uma divisão fracionária?

\begin{enumerate}
  \item \mintinline{haskell}{show (length xs) / 2}
  \item \mintinline{haskell}{fromRational (length xs) / 2}
  \item \mintinline{haskell}{read (length xs) / 2}
  \item \mintinline{haskell}{fromIntegral (length xs) / 2}
  \item \mintinline{haskell}{length xs `div` 2.0}

\end{enumerate}

\questionTitle{24}{Escolha entre div e barra}
\begin{flushright}
  \footnotesize
  \textit{\textbf{7: Polimorfismo}}\\
  \textit{cap07-expressao-div-vs-barra}\\

\end{flushright}
Qual alternativa compara corretamente \haskellinline{div} e \haskellinline{(/)}?
\begin{enumerate}
  \item \haskellinline{div} é uma função de \haskellinline{Show}, enquanto \haskellinline{(/)} é uma função de \haskellinline{Read}.
  \item As duas funções têm exatamente a mesma assinatura e diferem apenas no símbolo usado.
  \item \haskellinline{div} exige tipos \haskellinline{Integral}; \haskellinline{(/)} exige tipos \haskellinline{Fractional}.
  \item \haskellinline{div} é usada para \haskellinline{Double}; \haskellinline{(/)} é usada para \haskellinline{Int}.
  \item \haskellinline{(/)} sempre retorna \haskellinline{Int}, mesmo quando recebe \haskellinline{Double}.

\end{enumerate}

\questionTitle{25}{Enumeração com passo e consumo finito}
\begin{flushright}
  \footnotesize
  \textit{\textbf{6: Estruturas de dados básicas}}\\
  \textit{c06-listas-operacoes-03}\\

\end{flushright}
Considere a expressão:
\begin{minted}{haskell}
take 4 [3, 3 + 3 ..]
\end{minted}
Qual é o seu resultado?

\begin{enumerate}
  \item \haskellinline{ [3,6,9] }

  \item \haskellinline{ [3,6,9,12] }

  \item \haskellinline{ [6,9,12,15] }

  \item A expressão produz a lista infinita completa antes de aplicar \haskellinline{take}.

  \item \haskellinline{ [3,6,12,24] }


\end{enumerate}

\questionTitle{26}{Associatividade de subtração e potenciação}
\begin{flushright}
  \footnotesize
  \textit{\textbf{3: Expressões e definições}}\\
  \textit{c03-precedencia-associatividade-03}\\

\end{flushright}
Sabendo que \haskellinline|-| associa à esquerda e \haskellinline|^| associa à direita, qual alternativa está correta?

\begin{enumerate}
  \item Tanto \haskellinline|-| quanto \haskellinline|^| associam sempre à esquerda.

  \item \haskellinline|10 - 5 - 2| é lida como \haskellinline|(10 - 5) - 2|, enquanto \haskellinline|2 ^ 3 ^ 2| é lida como \haskellinline|2 ^ (3 ^ 2)|.

  \item Tanto \haskellinline|-| quanto \haskellinline|^| associam sempre à direita.

  \item A associatividade só se aplica a funções prefixas, não a operadores infixos.

  \item \haskellinline|10 - 5 - 2| é lida como \haskellinline|10 - (5 - 2)|, enquanto \haskellinline|2 ^ 3 ^ 2| é lida como \haskellinline|(2 ^ 3) ^ 2|.


\end{enumerate}

\questionTitle{27}{Otherwise fora da última guarda}
\begin{flushright}
  \footnotesize
  \textit{\textbf{5: Expressão condicional}}\\
  \textit{c05-otherwise-03}\\

\end{flushright}
Considere a definição:
\begin{haskellcode}
g x
  | otherwise = 1
  | x > 0     = 2
\end{haskellcode}
O que acontece ao avaliar \haskellinline|g 5|?

\begin{enumerate}
  \item O compilador reordena as guardas e produz \haskellinline|2|.

  \item Ocorre um erro em tempo de execução por ambiguidade entre as guardas.

  \item O resultado é \haskellinline|2|, porque a guarda \haskellinline|x > 0| seria mais específica.

  \item O resultado é \haskellinline|1|, porque \haskellinline|otherwise| equivale a \haskellinline|True|.

  \item Ocorre um erro de sintaxe, pois \haskellinline|otherwise| só pode aparecer na última linha.


\end{enumerate}

\questionTitle{28}{Classificação por faixas}
\begin{flushright}
  \footnotesize
  \textit{\textbf{5: Expressão condicional}}\\
  \textit{c05-guardas-avaliacao-03}\\

\end{flushright}
Analise a definição:
\begin{haskellcode}
conceito n
  | n >= 9    = "A"
  | n >= 7    = "B"
  | n >= 5    = "C"
  | otherwise = "D"
\end{haskellcode}
Qual é o resultado de \haskellinline|conceito 7.5|?

\begin{enumerate}
  \item \haskellinline|"A"|

  \item \haskellinline|"B"|

  \item \haskellinline|"C"|

  \item \haskellinline|"D"|

  \item Ocorre um erro de tipos, pois as guardas usam números decimais.


\end{enumerate}

\questionTitle{29}{Regra para abreviação de comandos}
\begin{flushright}
  \footnotesize
  \textit{\textbf{2: Ambiente de desenvolvimento}}\\
  \textit{c02-ghci-inspecao-03}\\

\end{flushright}
No ambiente GHCi, os comandos que começam com dois-pontos (\texttt{:}) frequentemente podem ser abreviados (por exemplo, \texttt{:r} no lugar de \texttt{:reload}). Qual é a regra técnica exigida pelo ambiente para que uma abreviação seja aceita?

\begin{enumerate}
  \item Apenas comandos que controlam o fluxo do terminal (como sair ou carregar) suportam atalhos; comandos de inspeção devem ser sempre digitados por extenso.

  \item A abreviação informada deve ser não ambígua, ou seja, deve ser um prefixo exclusivo que não coincida com a abreviação de nenhum outro comando existente.

  \item A abreviação só será aceita se o programador tiver configurado previamente um arquivo de manifesto indicando quais atalhos ele deseja mapear.

  \item A abreviação deve ter exatamente duas letras para que o analisador do GHCi não as confunda com nomes de variáveis Haskell válidas.

  \item O GHCi aceita qualquer abreviação independentemente de colisão; em caso de empate, ele sempre executa o último comando digitado no histórico da sessão.


\end{enumerate}

\questionTitle{30}{Assinatura em expressões}
\begin{flushright}
  \footnotesize
  \textit{\textbf{4: Tipos de Dados}}\\
  \textit{c04-assinatura-expressao-01}\\

\end{flushright}
Por que a expressão \haskellinline{read "42"} sozinha causaria um erro de ambiguidade de tipo, enquanto \haskellinline{read "42" :: Int} é aceita?

\begin{enumerate}
  \item \haskellinline{read} é sobrecarregada e pode retornar vários tipos; a anotação \haskellinline{:: Int} resolve a ambiguidade.

  \item A anotação é necessária porque \haskellinline{read} é uma função parcial.

  \item A anotação \haskellinline{:: Int} instrui o compilador a ignorar a inferência de tipos.

  \item \haskellinline{read} só funciona com anotação de tipo; sem ela, ocorre um erro de sintaxe.

  \item A anotação \haskellinline{:: Int} força a conversão da string para inteiro em tempo de execução.


\end{enumerate}

\questionTitle{31}{Tipo unitário}
\begin{flushright}
  \footnotesize
  \textit{\textbf{6: Estruturas de dados básicas}}\\
  \textit{c06-tuplas-unitario-01}\\

\end{flushright}
Qual afirmação descreve corretamente o tipo \haskellinline{()} em Haskell?

\begin{enumerate}
  \item É apenas outra notação para a lista vazia \haskellinline{[]}.

  \item É a forma como Haskell representa uma tupla de um elemento.

  \item É o tipo unitário: possui exatamente um único valor, também escrito como \haskellinline{()}.

  \item É um sinônimo de tipo para \haskellinline{Bool}.

  \item É um tipo reservado exclusivamente para expressões numéricas.


\end{enumerate}

\questionTitle{32}{Chamada que se afasta do caso base}
\begin{flushright}
  \footnotesize
  \textit{\textbf{9: Funções recursivas}}\\
  \textit{c09-fundamentos-terminacao-02}\\

\end{flushright}
Considere a função:
\begin{haskellcode}
potDois' :: Integer -> Integer
potDois' n
  | n == 0    = 1
  | otherwise = 2 * potDois' (n - 1)
\end{haskellcode}
Qual é o problema ao avaliar \haskellinline|potDois' (-5)|?

\begin{enumerate}
  \item Ocorre erro de tipo, pois \haskellinline|Integer| não permite valores negativos.

  \item A função retorna \haskellinline|1|, pois todo expoente negativo é tratado como o caso base.

  \item A cada chamada o argumento fica menor, afastando-se de \haskellinline|0|; portanto, a recursão não termina.

  \item A função retorna \haskellinline|0|, pois potências de dois negativas são arredondadas para zero.

  \item A função retorna \haskellinline|-32|, pois Haskell preserva automaticamente o sinal do expoente.


\end{enumerate}

\questionTitle{33}{Tipo de resultado de lookup}
\begin{flushright}
  \footnotesize
  \textit{\textbf{6: Estruturas de dados básicas}}\\
  \textit{c06-lookup-associacao-03}\\

\end{flushright}
Ao buscar uma chave em uma lista de associação com \haskellinline{lookup}, por que o resultado é representado com \haskellinline{Maybe}?

\begin{enumerate}
  \item Porque o uso de \haskellinline{lookup} transforma automaticamente qualquer lista em uma enumeração.

  \item Porque \haskellinline{lookup} sempre devolve uma lista, mesmo quando encontra apenas um valor.

  \item Porque Haskell não permite devolver diretamente valores como \haskellinline{String} ou \haskellinline{Double}.

  \item Porque toda função que opera sobre tuplas em Haskell precisa devolver um \haskellinline{Maybe}.

  \item Porque a chave pode estar presente ou ausente, e \haskellinline{Maybe} representa explicitamente essas duas possibilidades.


\end{enumerate}

\questionTitle{34}{A inovação do Sistema de Tipos da ML}
\begin{flushright}
  \footnotesize
  \textit{\textbf{1: Paradigmas de programação}}\\
  \textit{c01-historia-evolucao-03}\\

\end{flushright}
Na década de 1970, o pesquisador Robin Milner e sua equipe desenvolveram a linguagem ML (\emph{MetaLanguage}), projetada inicialmente para apoiar a demonstração automática de teoremas. A linguagem ML legou um avanço formidável, hoje presente fortemente em Haskell e linguagens como Rust e Swift. Qual foi essa inovação técnica?

\begin{enumerate}
  \item A invenção pioneira da \textbf{Orientação a Objetos}, que permitiu à indústria mesclar funções puras com o encapsulamento hermético de estado dentro de classes herdáveis.

  \item A arquitetura inovadora do compilador \textbf{JIT} (\emph{Just-in-Time Compiler}), que otimizou o código funcional para superar a velocidade dos programas escritos em linguagem C de forma absoluta.

  \item O mecanismo de \textbf{Avaliação Preguiçosa} estrita, que impedia o cálculo precipitado de matrizes infinitas e protegia os \emph{mainframes} contra falhas de estouro de memória.

  \item Um sistema de \textbf{inferência de tipos} elegante e poderoso, capaz de garantir a segurança estática sem exigir que o programador preenchesse o código com anotações manuais redundantes de tipos.

  \item A introdução e formalização das \textbf{Mônadas} na biblioteca padrão como a única abstração autorizada pelo compilador para isolar e manipular efeitos de entrada e saída.


\end{enumerate}

\questionTitle{35}{O papel de hFlush}
\begin{flushright}
  \footnotesize
  \textit{\textbf{8: Programas interativos}}\\
  \textit{c08-buffering-02}\\

\end{flushright}
Qual ação pode ser usada imediatamente depois de \haskellinline|putStr| para garantir que o conteúdo pendente na saída padrão seja enviado ao terminal?

\begin{enumerate}
  \item \haskellinline|delay stdout|.

  \item \haskellinline|hClear stdout|.

  \item \haskellinline|hRefresh stdout|.

  \item \haskellinline|hFlush stdout|.

  \item \haskellinline|getLine stdout|.


\end{enumerate}

\questionTitle{36}{O papel de main}
\begin{flushright}
  \footnotesize
  \textit{\textbf{8: Programas interativos}}\\
  \textit{c08-modelo-io-02}\\

\end{flushright}
Qual é o papel especial de uma definição como \haskellinline|main :: IO ()| em um programa Haskell executável?

\begin{enumerate}
  \item \haskellinline|main| é apenas uma convenção de nomenclatura; qualquer função do arquivo poderia ser executada automaticamente.

  \item \haskellinline|main| é a única definição do programa que pode chamar funções puras.

  \item \haskellinline|main| é obrigatoriamente uma função pura, pois programas Haskell não podem ter efeitos de E/S.

  \item \haskellinline|main| é a ação de E/S que o \emph{runtime system} executa quando o programa é iniciado.

  \item \haskellinline|main| deve retornar um código inteiro de saída, como ocorre em C.


\end{enumerate}

\questionTitle{37}{Projeto para cálculo de média}
\begin{flushright}
  \footnotesize
  \textit{\textbf{10: Ações de E/S recursivas}}\\
  \textit{c10-separacao-puro-03}\\

\end{flushright}
Para o exercício da média de uma sequência, considere o objetivo: ler valores até um sentinela negativo e imprimir a média dos valores não negativos lidos. Qual decomposição é mais alinhada ao estilo do capítulo?

\begin{enumerate}
  \item Uma definição que ignora a lista vazia, pois \haskellinline|fromIntegral (length lista)| nunca pode ser zero.

  \item Uma ação \haskellinline|lerLista :: IO [Double]| para coletar os valores e uma função pura ou expressão separada para calcular a média da lista, tratando a lista vazia antes da divisão.

  \item Uma função \haskellinline|sum :: IO [Double] -> Double|, pois \haskellinline|sum| opera diretamente sobre ações de E/S.

  \item Uma ação \haskellinline|main| que acumula uma string com todos os números e chama \haskellinline|read| apenas no final.

  \item Uma única função pura \haskellinline|media :: IO Double| que chama \haskellinline|readLn| internamente sem usar \haskellinline|do|.


\end{enumerate}

\questionTitle{38}{A analogia da receita de bolo}
\begin{flushright}
  \footnotesize
  \textit{\textbf{1: Paradigmas de programação}}\\
  \textit{c01-conceitos-paradigmas-03}\\

\end{flushright}
Ao explicar paradigmas de programação, é comum utilizar a analogia de uma \emph{receita de bolo} (ex: \emph{1. quebre os ovos; 2. adicione farinha; 3. asse por 40 minutos}). 

Esta analogia ilustra de forma primária o raciocínio de qual paradigma e por quê?

\begin{enumerate}
  \item O paradigma orientado a objetos, pois a receita representa a troca de mensagens assíncronas entre as instâncias dos ingredientes.

  \item O paradigma funcional puro, pois cada passo da receita é uma função referencialmente transparente que não altera o estado da cozinha.

  \item O paradigma declarativo, pois descreve as propriedades finais do bolo sem se preocupar com a ordem matemática em que os ingredientes foram misturados.

  \item O paradigma imperativo, pois foca na descrição passo a passo da sequência de controle e nas mudanças sucessivas de estado para atingir o objetivo.

  \item O paradigma lógico, pois cada passo da receita atua como um axioma em um motor de inferência que deduz a existência do bolo.


\end{enumerate}

\questionTitle{39}{Projeto de fatorial duplo}
\begin{flushright}
  \footnotesize
  \textit{\textbf{9: Funções recursivas}}\\
  \textit{c09-projeto-fatorial-duplo-01}\\

\end{flushright}
Qual definição expressa corretamente o fatorial duplo de um número natural, reduzindo o argumento de dois em dois?

\begin{enumerate}
  \item \begin{haskellcode}
fatorialDuplo n
  | n == 0 = 1
  | n > 0  = fatorialDuplo (n + 2)
\end{haskellcode}

  \item \begin{haskellcode}
fatorialDuplo n
  | n == 0 = 1
  | n == 1 = 1
  | n > 1  = n * fatorialDuplo (n - 2)
\end{haskellcode}

  \item \begin{haskellcode}
fatorialDuplo n
  | n == 0 = 0
  | n > 0  = n * fatorialDuplo (n - 2)
\end{haskellcode}

  \item \begin{haskellcode}
fatorialDuplo n
  | n == 1 = 1
  | n > 1  = n * fatorialDuplo (n - 1)
\end{haskellcode}

  \item \begin{haskellcode}
fatorialDuplo n
  | n == 0 = 1
  | n > 0  = n + fatorialDuplo (n - 2)
\end{haskellcode}


\end{enumerate}

\questionTitle{40}{Papel do acumulador}
\begin{flushright}
  \footnotesize
  \textit{\textbf{9: Funções recursivas}}\\
  \textit{c09-cauda-acumulador-01}\\

\end{flushright}
Em uma função recursiva com acumulador, qual é o papel típico do parâmetro acumulador?

\begin{enumerate}
  \item Impedir que a função seja chamada com números negativos.

  \item Guardar o resultado parcial já construído, permitindo que a chamada recursiva avance levando esse estado adiante.

  \item Armazenar automaticamente todos os argumentos anteriores em uma lista.

  \item Substituir o caso base, tornando-o desnecessário.

  \item Fazer com que o compilador transforme qualquer função em uma função não recursiva.


\end{enumerate}

\questionTitle{41}{Avaliação lazy e argumento não usado}
\begin{flushright}
  \footnotesize
  \textit{\textbf{3: Expressões e definições}}\\
  \textit{c03-layout-avaliacao-03}\\

\end{flushright}
Considere a expressão:

\begin{haskellcode}
(\x -> 1) (div 1 0)
\end{haskellcode}

O que acontece em Haskell, considerando a avaliação \emph{lazy}?

\begin{enumerate}
  \item A expressão avalia para \haskellinline|0|, pois Haskell define \haskellinline|div 1 0| como zero em contextos não usados.

  \item A expressão avalia para \haskellinline|1|, pois o argumento \haskellinline|div 1 0| não é necessário para produzir o resultado da função.

  \item A expressão avalia para a própria função \haskellinline|\x -> 1|, sem aplicar o argumento.

  \item A expressão sempre produz erro de divisão por zero antes de a função ser aplicada, como ocorreria em avaliação estrita.

  \item A expressão é sintaticamente inválida, pois uma lambda não pode ser aplicada diretamente a um argumento.


\end{enumerate}

\questionTitle{42}{Inferência de tipo de função composta}
\begin{flushright}
  \footnotesize
  \textit{\textbf{4: Tipos de Dados}}\\
  \textit{c04-inferencia-funcao-01}\\

\end{flushright}
Considere a função definida por \haskellinline{f x = not (x > 0)}. Qual é o tipo inferido para \haskellinline{f}?

\begin{enumerate}
  \item \haskellinline{Double -> Bool}

  \item \haskellinline{Int -> Bool}

  \item \haskellinline{a -> Bool}

  \item \haskellinline{Num a => a -> Bool}

  \item \haskellinline{(Ord a, Num a) => a -> Bool}


\end{enumerate}

\questionTitle{43}{Let como subexpressão}
\begin{flushright}
  \footnotesize
  \textit{\textbf{3: Expressões e definições}}\\
  \textit{c03-let-escopo-03}\\

\end{flushright}
Qual é o valor da expressão \haskellinline|2 * (let x = 3; y = 4 in x + y)|?

\begin{enumerate}
  \item \haskellinline|14|.

  \item A expressão é inválida, pois \haskellinline|let| não pode aparecer dentro de parênteses.

  \item \haskellinline|24|.

  \item \haskellinline|7|.

  \item \haskellinline|9|.


\end{enumerate}

\questionTitle{44}{sort como combinação de estrutura paramétrica e restrição}
\begin{flushright}
  \footnotesize
  \textit{\textbf{7: Polimorfismo}}\\
  \textit{cap07-comparacao-sort-param-adhoc}\\

\end{flushright}
A assinatura abaixo combina duas ideias:

\begin{minted}{haskell}
sort :: Ord a => [a] -> [a]
\end{minted}

Qual interpretação é mais adequada?

\begin{enumerate}
  \item A função é uma conversão numérica, pois \haskellinline{Ord} pertence apenas à hierarquia de \haskellinline{Num}.
  \item A estrutura da função é polimórfica sobre listas de \haskellinline{a}, mas a ordenação dos elementos exige uma instância \haskellinline{Ord a}.
  \item A presença de \haskellinline{Ord a} torna a função monomórfica, aceitando apenas \haskellinline{Int}.
  \item A função depende de herança entre objetos mutáveis.
  \item A assinatura mostra polimorfismo paramétrico puro, pois classes de tipo não impõem restrições reais.

\end{enumerate}

\questionTitle{45}{Paralelismo Seguro}
\begin{flushright}
  \footnotesize
  \textit{\textbf{1: Paradigmas de programação}}\\
  \textit{c01-engenharia-software-03}\\

\end{flushright}
No contexto de sistemas de alta disponibilidade executados em processadores \emph{multicore}, a programação funcional tem ganhado vasto espaço na indústria contemporânea. Qual característica fundacional do paradigma facilita imensamente a construção de programas paralelos livres de falhas?

\begin{enumerate}
  \item O paradigma impossibilita deliberadamente a criação de processos paralelos (sendo restrito a \emph{single-thread}), o que mitiga os erros de concorrência anulando a complexidade do escalonamento.

  \item A imutabilidade dos dados combinada às funções puras garantem que múltiplas \emph{threads} não causarão \emph{race conditions} (condições de corrida) tentando modificar o mesmo espaço de memória simultaneamente.

  \item Os compiladores de linguagens funcionais traduzem nativamente o código para chamadas assíncronas de \emph{hardware} que contornam totalmente os barramentos lentos da memória principal.

  \item A avaliação estrita das linguagens funcionais impõe um controle rígido ao escalonador do sistema operacional, forçando a distribuição milimétrica da carga da aplicação entre os núcleos físicos.

  \item A tipagem dinâmica fraca da linguagem assegura que os processos em segundo plano reescrevam as estruturas de memória de forma atômica e sincronizada por meio do \emph{garbage collector}.


\end{enumerate}

\questionTitle{46}{Inferência de literais no GHCi (parametrizado)}
\begin{flushright}
  \footnotesize
  \textit{\textbf{4: Tipos de Dados}}\\
  \textit{c04-ghci-inferencia-01}\\

\end{flushright}
Qual tipo o GHCi inferiria para o literal \haskellinline{'x' } ao usar o comando \texttt{:type}?

\begin{enumerate}
  \item \haskellinline{String}

  \item Ocorre um erro de sintaxe.

  \item \haskellinline{Letter}

  \item \haskellinline{Char}

  \item \haskellinline{[Char]}


\end{enumerate}

\questionTitle{47}{Inferência com aplicação de função e tuplas}
\begin{flushright}
  \footnotesize
  \textit{\textbf{7: Polimorfismo}}\\
  \textit{cap07-inferencia-tupla-funcoes}\\

\end{flushright}
Considere:

\begin{minted}{haskell}
aplicaGuarda f x = (f x, x)
\end{minted}

Qual é a assinatura de tipo mais geral?

\begin{enumerate}
  \item \mintinline{haskell}{aplicaGuarda :: (a -> a) -> a -> (a, a)}
  \item \mintinline{haskell}{aplicaGuarda :: Eq a => (a -> b) -> a -> (b, a)}
  \item \mintinline{haskell}{aplicaGuarda :: (a -> b) -> a -> (b, a)}
  \item \mintinline{haskell}{aplicaGuarda :: [a] -> a -> (a, a)}
  \item \mintinline{haskell}{aplicaGuarda :: (a -> b) -> b -> (b, a)}

\end{enumerate}

\questionTitle{48}{Identificando tipos vs. variáveis}
\begin{flushright}
  \footnotesize
  \textit{\textbf{4: Tipos de Dados}}\\
  \textit{c04-conceito-nomes-02}\\

\end{flushright}
Considere os identificadores: \haskellinline{Pessoa} e \haskellinline{pessoa}. Qual é a interpretação mais provável em Haskell?

\begin{enumerate}
  \item \haskellinline{Pessoa} é um tipo (ou construtor de dados) e \haskellinline{pessoa} é uma variável ou função.

  \item Ambos são tipos, sendo \haskellinline{pessoa} um sinônimo.

  \item \haskellinline{Pessoa} é uma variável global e \haskellinline{pessoa} é uma variável local.

  \item Não há diferença; Haskell não distingue maiúsculas de minúsculas.

  \item Ambos são variáveis, mas com escopos diferentes.


\end{enumerate}

\questionTitle{49}{Lista final produzida por acumulador com reverse}
\begin{flushright}
  \footnotesize
  \textit{\textbf{10: Ações de E/S recursivas}}\\
  \textit{c10-listas-ordem-03}\\

\end{flushright}
A função abaixo é executada inicialmente como \haskellinline|lerLista []|. O usuário digita \haskellinline|4, 4, 3, 2, 0|, um valor por linha.

\begin{haskellcode}
lerLista :: [Integer] -> IO [Integer]
lerLista xs = do
  x <- readLn
  if x == 0
    then return (reverse xs)
    else lerLista (x:xs)
\end{haskellcode}

Qual lista é produzida pela ação?

\begin{enumerate}
  \item \haskellinline|[4,4,3,2]|

  \item \haskellinline|[4,4,3,2,0]|

  \item \haskellinline|[0]|

  \item \haskellinline|[2,3,4,4]|

  \item \haskellinline|[]|


\end{enumerate}

\questionTitle{50}{Soma de dois números}
\begin{flushright}
  \footnotesize
  \textit{\textbf{8: Programas interativos}}\\
  \textit{c08-programas-completos-02}\\

\end{flushright}
Considere:

\begin{minted}{haskell}
main :: IO ()
main = do
  x <- readLn :: IO Int
  y <- readLn :: IO Int
  print (x + y)
\end{minted}

Se o usuário digitar \haskellinline|3| e depois \haskellinline|4|, qual será a saída produzida por \haskellinline|print|?

\begin{enumerate}
  \item \haskellinline|(3,4)|

  \item \haskellinline|3 + 4|

  \item \haskellinline|7|

  \item Nenhuma saída, pois \haskellinline|readLn| não pode ser usado dentro de \haskellinline|do|.

  \item \haskellinline|"34"|


\end{enumerate}

\questionTitle{51}{Uso seguro de fromJust}
\begin{flushright}
  \footnotesize
  \textit{\textbf{6: Estruturas de dados básicas}}\\
  \textit{c06-maybe-operacoes-02}\\

\end{flushright}
Em qual situação o uso de \haskellinline{fromJust} é mais justificável neste capítulo?

\begin{enumerate}
  \item Sempre que a expressão tiver o tipo \haskellinline{Maybe Int}.

  \item Sempre que a expressão aparecer dentro de uma lista.

  \item Sempre que quisermos evitar o uso de \haskellinline{fromMaybe}.

  \item Apenas quando o conteúdo for uma string, mas não quando for um número.

  \item Quando uma verificação anterior, como \haskellinline{isJust}, já garantiu que o valor não é \haskellinline{Nothing}.


\end{enumerate}

\questionTitle{52}{Falha de conversão com read}
\begin{flushright}
  \footnotesize
  \textit{\textbf{8: Programas interativos}}\\
  \textit{c08-conversao-validacao-02}\\

\end{flushright}
Considere o programa:

\begin{minted}{haskell}
main :: IO ()
main = do
  putStrLn "Digite um número:"
  linha <- getLine
  let numero = read linha :: Int
  print (numero * 2)
\end{minted}

O que acontece se o usuário digitar \haskellinline|"cinco"|?

\begin{enumerate}
  \item O programa não compila, pois o compilador detecta estaticamente que o usuário pode digitar texto inválido.

  \item O programa imprime \haskellinline|"cincocinco"|, pois o operador \haskellinline|*| repete strings.

  \item O programa chama \haskellinline|getLine| novamente até que o usuário digite um número.

  \item A conversão com \haskellinline|read| falha e o programa termina com uma exceção em tempo de execução quando \haskellinline|numero| é avaliado.

  \item O programa imprime \haskellinline|0|, pois \haskellinline|read| retorna zero quando a conversão falha.


\end{enumerate}

\questionTitle{53}{Representação de caracteres}
\begin{flushright}
  \footnotesize
  \textit{\textbf{4: Tipos de Dados}}\\
  \textit{c04-tipos-char-01}\\

\end{flushright}
Qual das seguintes alternativas representa o caractere de nova linha (newline) em Haskell?

\begin{enumerate}
  \item \haskellinline{"\n"}

  \item \haskellinline{'\n'}

  \item \haskellinline{'\newline'}

  \item \haskellinline{'\\n'}

  \item \haskellinline{"\LF"}


\end{enumerate}

\questionTitle{54}{print e putStrLn}
\begin{flushright}
  \footnotesize
  \textit{\textbf{8: Programas interativos}}\\
  \textit{c08-entrada-saida-02}\\

\end{flushright}
Qual é a diferença central entre \haskellinline|print| e \haskellinline|putStrLn|?

\begin{enumerate}
  \item \haskellinline|print| lê um valor da entrada; \haskellinline|putStrLn| imprime um valor na saída.

  \item \haskellinline|print| aplica \haskellinline|show| ao valor antes de imprimir; \haskellinline|putStrLn| recebe diretamente uma \haskellinline|String| e acrescenta uma quebra de linha.

  \item Não há diferença; as duas funções têm exatamente o mesmo tipo e o mesmo comportamento.

  \item \haskellinline|print| só funciona para \haskellinline|String|; \haskellinline|putStrLn| funciona para qualquer tipo da classe \haskellinline|Show|.

  \item \haskellinline|print| imprime sem quebra de linha; \haskellinline|putStrLn| imprime com quebra de linha.


\end{enumerate}

\questionTitle{55}{If produzindo função}
\begin{flushright}
  \footnotesize
  \textit{\textbf{5: Expressão condicional}}\\
  \textit{c05-if-avaliacao-03}\\

\end{flushright}
Qual é o resultado da expressão abaixo?
\begin{haskellcode}
(if 5 > 3 then (*) else (+)) 4 2
\end{haskellcode}

\begin{enumerate}
  \item \haskellinline|14|

  \item Ocorre um erro de sintaxe, pois \haskellinline|if| não pode aparecer entre parênteses.

  \item \haskellinline|8|

  \item Ocorre um erro de tipo, pois os ramos do \haskellinline|if| são funções.

  \item \haskellinline|6|


\end{enumerate}

\questionTitle{56}{Segurança de tipos com sinônimos}
\begin{flushright}
  \footnotesize
  \textit{\textbf{4: Tipos de Dados}}\\
  \textit{c04-sinonimos-seguranca-01}\\

\end{flushright}
Considere as definições \haskellinline{type UsuarioID = Int} e \haskellinline{type ProdutoID = Int}. Se uma função espera um \haskellinline{UsuarioID} mas recebe um \haskellinline{ProdutoID}, o que acontecerá?

\begin{enumerate}
  \item O programa compilará, mas emitirá um aviso (\emph{warning}).

  \item A palavra-chave \haskellinline{type} não pode ser usada com tipos numéricos.

  \item O programa compilará normalmente, pois para o sistema de tipos, \haskellinline{UsuarioID}, \haskellinline{ProdutoID} e \haskellinline{Int} são o mesmo tipo.

  \item O programa falhará na compilação com um erro de incompatibilidade de tipos.

  \item Ocorrerá um erro em tempo de execução.


\end{enumerate}

\questionTitle{57}{Escolha entre Int e Integer}
\begin{flushright}
  \footnotesize
  \textit{\textbf{9: Funções recursivas}}\\
  \textit{c09-eficiencia-integer-int-01}\\

\end{flushright}
Por que o capítulo usa \haskellinline|Integer|, e não \haskellinline|Int|, em exemplos como \haskellinline|fatorial|?

\begin{enumerate}
  \item Porque \haskellinline|Int| só pode ser usado em listas.

  \item Porque \haskellinline|Integer| torna qualquer função automaticamente recursiva de cauda.

  \item Porque \haskellinline|Integer| tem precisão arbitrária, enquanto \haskellinline|Int| tem tamanho fixo e pode sofrer estouro de capacidade.

  \item Porque \haskellinline|Integer| impede chamadas recursivas infinitas.

  \item Porque \haskellinline|Int| não permite valores positivos.


\end{enumerate}

\questionTitle{58}{Ordem das definições locais em where}
\begin{flushright}
  \footnotesize
  \textit{\textbf{3: Expressões e definições}}\\
  \textit{c03-definicoes-where-03}\\

\end{flushright}
Considere a definição:

\begin{haskellcode}
g x = a * b
  where
    b = x + 2
    a = b - 1
\end{haskellcode}

Qual é o valor de \haskellinline|g 4|?

\begin{enumerate}
  \item \haskellinline|20|.

  \item Ocorre erro, pois definições em \haskellinline|where| não podem depender umas das outras.

  \item \haskellinline|18|.

  \item Ocorre erro, pois \haskellinline|a| aparece antes de \haskellinline|b| no corpo da função.

  \item \haskellinline|30|.


\end{enumerate}

\questionTitle{59}{Quantidade de leituras até o sentinela}
\begin{flushright}
  \footnotesize
  \textit{\textbf{10: Ações de E/S recursivas}}\\
  \textit{c10-rastreamento-02}\\

\end{flushright}
Considere a mesma função:

\begin{haskellcode}
lerESomar :: IO Integer
lerESomar = do
  n <- readLn
  if n == 0
    then return 0
    else do
      somaResto <- lerESomar
      return (n + somaResto)
\end{haskellcode}

Se o usuário digita \haskellinline|6, 5, 4, 3, 0|, quantas chamadas a \haskellinline|readLn| são realizadas?

\begin{enumerate}
  \item 5

  \item 18

  \item 4

  \item 0

  \item 1


\end{enumerate}

\questionTitle{60}{Assinatura de função com Char, Int e String}
\begin{flushright}
  \footnotesize
  \textit{\textbf{4: Tipos de Dados}}\\
  \textit{c04-funcao-assinatura-03}\\

\end{flushright}
Qual é a assinatura de tipo de uma função que recebe um \haskellinline{Char}, um \haskellinline{Int} e retorna uma \haskellinline{String}?

\begin{enumerate}
  \item \haskellinline{(Char, Int) -> String}

  \item \haskellinline{[Char, Int, String]}

  \item \haskellinline{Char -> Int -> String}

  \item \haskellinline{Char, Int -> String}

  \item \haskellinline{String -> Char -> Int}


\end{enumerate}

\questionTitle{61}{Conflito de classes em checagem estática}
\begin{flushright}
  \footnotesize
  \textit{\textbf{4: Tipos de Dados}}\\
  \textit{c04-checagem-conflito-01}\\

\end{flushright}
Considere o código \haskellinline{calculo x y = (x / 2) `div` y}. Por que este código falha na compilação?

\begin{enumerate}
  \item A notação infixa com crases não pode ser usada com operadores simbólicos.

  \item A função \haskellinline{div} tem maior precedência que \haskellinline{/}.

  \item O operador \haskellinline{/} exige um tipo fracionário (\haskellinline{Fractional}), mas \haskellinline{div} exige um tipo integral (\haskellinline{Integral}), causando um conflito de tipos.

  \item O literal \haskellinline{2} é inferido como \haskellinline{Integer}, que não é compatível com \haskellinline{x}.

  \item Falta uma anotação de tipo explícita, impedindo a inferência.


\end{enumerate}

\questionTitle{62}{Papel de where em guardas}
\begin{flushright}
  \footnotesize
  \textit{\textbf{5: Expressão condicional}}\\
  \textit{c05-where-escopo-03}\\

\end{flushright}
Qual é a principal vantagem de usar \haskellinline|where| em uma definição com guardas, como no cálculo de \haskellinline|delta| em uma função que determina o número de raízes de uma equação?

\begin{enumerate}
  \item Garantir que todas as guardas sejam exaustivas.

  \item Tornar a função automaticamente polimórfica.

  \item Permitir que a função tenha ramos com tipos diferentes.

  \item Permitir nomear um cálculo local e reutilizá-lo em guardas e resultados, evitando repetição.

  \item Fazer com que as guardas sejam avaliadas em paralelo.


\end{enumerate}

\questionTitle{63}{Modelagem composta de pedidos}
\begin{flushright}
  \footnotesize
  \textit{\textbf{6: Estruturas de dados básicas}}\\
  \textit{c06-modelagem-03}\\

\end{flushright}
Qual tipo representa melhor uma coleção de pedidos em que cada pedido possui o nome do cliente, a lista de itens comprados e um total opcional?

\begin{enumerate}
  \item \haskellinline{[(String, [(String, Int)], Maybe Double)]}

  \item \haskellinline{[String, [(String, Int)], Maybe Double]}

  \item \haskellinline{([String], [Int], [Double])}

  \item \haskellinline{(String, [String], Double)}

  \item \haskellinline{Maybe [(String, Int, Double)]}


\end{enumerate}

\questionTitle{64}{Erro comum com let}
\begin{flushright}
  \footnotesize
  \textit{\textbf{8: Programas interativos}}\\
  \textit{c08-blocos-do-02}\\

\end{flushright}
Um programador escreveu:

\begin{minted}{haskell}
main :: IO ()
main = do
  putStrLn "Qual é o seu nome?"
  let nome = getLine
  putStrLn ("Olá, " ++ nome)
\end{minted}

Qual é a causa fundamental do erro?

\begin{enumerate}
  \item A função \haskellinline|putStrLn| não aceita expressões com concatenação de strings.

  \item \haskellinline|nome| fica associado à ação \haskellinline|getLine|, de tipo \haskellinline|IO String|. O operador \haskellinline|++| espera uma \haskellinline|String|, não uma ação \haskellinline|IO String|.

  \item A palavra \haskellinline|let| é sempre proibida dentro de blocos \haskellinline|do|.

  \item \haskellinline|getLine| não pode ser usado depois de \haskellinline|putStrLn| no mesmo bloco \haskellinline|do|.

  \item A variável \haskellinline|nome| precisava ser declarada globalmente antes de \haskellinline|main|.


\end{enumerate}

\questionTitle{65}{Cons versus concatenação}
\begin{flushright}
  \footnotesize
  \textit{\textbf{6: Estruturas de dados básicas}}\\
  \textit{c06-listas-estrutura-03}\\

\end{flushright}
Qual afirmação descreve corretamente a diferença entre \haskellinline{(:)} e \haskellinline{(++)}?

\begin{enumerate}
  \item \haskellinline{(:)} concatena listas mais rapidamente do que \haskellinline{(++)}, mas exige listas do mesmo tamanho.

  \item \haskellinline{(:)} adiciona um elemento ao início de uma lista, enquanto \haskellinline{(++)} concatena duas listas.

  \item \haskellinline{(:)} e \haskellinline{(++)} sempre produzem exatamente o mesmo efeito, mudando apenas a sintaxe.

  \item \haskellinline{(:)} adiciona um elemento ao final da lista, enquanto \haskellinline{(++)} adiciona ao início.

  \item \haskellinline{(++)} adiciona um elemento à lista, enquanto \haskellinline{(:)} exige duas listas como argumentos.


\end{enumerate}

\questionTitle{66}{Análise de IMC com where}
\begin{flushright}
  \footnotesize
  \textit{\textbf{5: Expressão condicional}}\\
  \textit{c05-where-avaliacao-03}\\

\end{flushright}
Considere a função:
\begin{haskellcode}
analisaIMC peso altura
  | imc < 18.5 = "Abaixo"
  | imc < 25.0 = "Normal"
  | otherwise  = "Acima"
  where
    imc = peso / (altura * altura)
\end{haskellcode}
Qual é o resultado de \haskellinline|analisaIMC 80 2.0|?

\begin{enumerate}
  \item \haskellinline|"Normal"|

  \item \haskellinline|"80.0"|

  \item Ocorre um erro, pois variáveis definidas em \haskellinline|where| não podem aparecer nas guardas.

  \item \haskellinline|"Abaixo"|

  \item \haskellinline|"Acima"|


\end{enumerate}

\questionTitle{67}{A distinção entre GHC e GHCi}
\begin{flushright}
  \footnotesize
  \textit{\textbf{2: Ambiente de desenvolvimento}}\\
  \textit{c02-ecossistema-ferramentas-03}\\

\end{flushright}
O conjunto de ferramentas do Glasgow Haskell Compiler fornece dois componentes essenciais: o \texttt{ghc} e o \texttt{ghci}. Qual é a distinção primária entre eles durante o desenvolvimento?

\begin{enumerate}
  \item O \texttt{ghc} traduz funções puras matemáticas, enquanto o \texttt{ghci} é reservado apenas para compilar funções que realizam operações de entrada e saída (I/O).

  \item Não há diferença funcional; ambos são atalhos distintos de terminal que iniciam exatamente a mesma ferramenta de \emph{build}.

  \item O \texttt{ghc} é o compilador de linha de comando que gera binários executáveis, enquanto o \texttt{ghci} é o ambiente interativo (\emph{REPL}) para avaliação em tempo real.

  \item O \texttt{ghc} é utilizado exclusivamente no sistema operacional Linux, enquanto o \texttt{ghci} é a versão portada para o Windows.

  \item O \texttt{ghc} gerencia a formatação do código no editor de texto, ao passo que o \texttt{ghci} executa os testes de unidade automatizados.


\end{enumerate}

\questionTitle{68}{Comando para recarregar um arquivo no GHCi}
\begin{flushright}
  \footnotesize
  \textit{\textbf{3: Expressões e definições}}\\
  \textit{c03-ghci-definicoes-03}\\

\end{flushright}
Após carregar um arquivo Haskell no GHCi e editar o código no editor, qual comando recarrega o último arquivo carregado?

\begin{enumerate}
  \item \haskellinline|:reload|, ou sua abreviação \haskellinline|:r|.

  \item \haskellinline|:browse|, pois ele atualiza o arquivo fonte em disco.

  \item \haskellinline|:edit|, pois ele salva e recompila automaticamente o módulo.

  \item \haskellinline|:load| sem argumento, pois ele sempre recarrega o arquivo anterior.

  \item \haskellinline|:type|, pois ele recompila e mostra os tipos de todas as definições.


\end{enumerate}

\questionTitle{69}{Valores e forma normal}
\begin{flushright}
  \footnotesize
  \textit{\textbf{3: Expressões e definições}}\\
  \textit{c03-comentarios-valores-03}\\

\end{flushright}
No contexto da avaliação de expressões em Haskell, qual alternativa caracteriza corretamente a relação entre expressões e valores?

\begin{enumerate}
  \item Uma expressão em forma normal ainda precisa ser reduzida pelo menos uma vez para produzir um valor.

  \item Todo valor é uma expressão, mas nem toda expressão já é um valor; uma expressão torna-se um valor quando atinge uma forma que não pode mais ser reduzida.

  \item Valores são apenas números; listas, tuplas e funções não são considerados valores em Haskell.

  \item Uma expressão só é considerada valor se tiver sido associada a um nome por uma definição.

  \item Toda expressão escrita pelo programador já é imediatamente um valor, mesmo antes de qualquer avaliação.


\end{enumerate}


\end{multicols}
\makeanswersheet{UFOP}{BCC222: Programação Funcional}{José Romildo}{2026/1}{Primeira Prova}{02}{69}{25}{7}
\gabarito{02}{B,D,B,C,D,B,E,B,B,B,C,B,B,B,D,B,C,E,C,C,A,E,D,C,B,B,D,B,B,A,C,C,E,D,D,D,B,D,B,B,B,E,A,B,B,D,C,A,A,C,E,D,B,B,C,C,C,E,A,C,C,D,A,B,B,A,C,A,B}
\end{document}
